\section{Engine and System Architecture}\label{sec:arch}

The engine is designed so that a step costs what the algorithm needs and no more: $N$ decisions, plus one coupling row for each spin that flips. All spin and field state lives on chip exactly once, and the datapath contains no handshakes. Earlier versions of the engine (v1--v3 in Table~\ref{tab:evo}) are written in high-level synthesis (HLS), which outlines every pipelined loop into its own module and copies the register-resident state into each one (356\,k flip-flops for 100\,k bits of state); the loop-exit handshakes also become clock enables with 32--43\,k fanout, which limits timing at 250--300\,MHz. The engine described here is hand-written RTL, and its only HLS part is a small front end that moves data (Figure~\ref{fig:arch}).

\begin{figure*}[t]
\centering
\definecolor{figtext}{HTML}{374D5C}\definecolor{figblue}{HTML}{DAE8FC}\definecolor{figgreen}{HTML}{D5E8D4}%
\definecolor{figorange}{HTML}{FFE6CC}\definecolor{figpurple}{HTML}{E1D5E7}\definecolor{figred}{HTML}{F8CECC}\definecolor{figredline}{HTML}{B85450}\definecolor{figyellow}{HTML}{FFF2CC}\definecolor{figgray}{HTML}{F5F5F5}%
\resizebox{\textwidth}{!}{%
\begin{tikzpicture}[
  font=\fontfamily{phv}\selectfont\footnotesize, >=Latex, text=figtext,
  blk/.style={draw=figtext, rounded corners=2pt, align=center, inner sep=2.5pt, fill=white, line width=0.5pt},
  lab/.style={align=center, font=\fontfamily{phv}\selectfont\scriptsize},
  bus/.style={->, line width=0.7pt, draw=figtext},
  bc/.style={->, line width=0.9pt, draw=figredline}
]
\node[anchor=north west, font=\fontfamily{phv}\selectfont\footnotesize\bfseries] at (0,0.3) {(a) AMD Alveo V80};
\draw[figtext, line width=0.6pt, rounded corners=2pt, fill=white] (0,-0.12) rectangle (3.7,-6.65);
\foreach \s/\k in {2/0,1/1,0/2} {
  \pgfmathsetmacro{\yt}{-0.62-\k*1.6}
  \draw[figtext, dashed, line width=0.5pt, fill=figgray] (0.15,\yt) rectangle (3.55,\yt-1.42);
  \node[anchor=north west, lab] at (0.18,\yt) {SLR\s};
  \foreach \e in {0,1,2,3} {
    \pgfmathtruncatemacro{\n}{4*\s+\e}
    \node[blk, fill=figorange, minimum width=0.66cm, minimum height=0.62cm, lab] (E\n) at (0.62+\e*0.76,\yt-0.86) {\n};
  }
}
\node[blk, fill=figgreen, minimum width=3.4cm, minimum height=0.7cm] at (1.85,-6.05) {NoC $\cdot$ HBM $\cdot$ PCIe};
\draw[figtext, line width=1.2pt, rounded corners=2pt] (E7.south west) rectangle (E7.north east);
\draw[gray!55, line width=0.4pt] (E7.north east) -- (4.1,-0.12);
\draw[gray!55, line width=0.4pt] (E7.south east) -- (4.1,-6.65);
\node[anchor=north west, font=\fontfamily{phv}\selectfont\footnotesize\bfseries] at (4.1,0.3) {(b) One engine: 32 groups, 256 decision lanes, 8 flip extractors};
\draw[figtext, line width=0.6pt, rounded corners=2pt, fill=white] (4.1,-0.12) rectangle (15.15,-6.65);
\node[blk, fill=figgreen, text width=2.0cm] (fe) at (5.3,-1.5) {HLS front end\\{\scriptsize AXI4 to HBM, timer}};
\node[blk, fill=figpurple, text width=2.0cm] (ld) at (5.3,-2.8) {Loader\\{\scriptsize header, $J$, tables}};
\node[blk, fill=figpurple, text width=2.0cm] (tab) at (5.3,-4.1) {Step tables\\{\scriptsize $4T_t$, $q_t$, $\gamma_t$, constant}};
\node[blk, fill=figpurple, text width=2.0cm] (sc) at (5.3,-5.4) {Score and output\\{\scriptsize $\sum_i s_ih_i$, cut, spins}};
\draw[bus] (fe) -- (ld);
\draw[bus] (ld) -- (tab);
\foreach \o in {0.28,0.14} \draw[figtext, line width=0.5pt, rounded corners=2pt, fill=figblue] (7.2+\o,-1.2+\o) rectangle (11.4+\o,-6.0+\o);
\draw[figtext, line width=0.6pt, rounded corners=2pt, fill=figblue] (7.2,-1.2) rectangle (11.4,-6.0);
\node[anchor=north west, font=\fontfamily{phv}\selectfont\scriptsize\bfseries] at (7.25,-1.23) {group $b$};
\draw[figtext, line width=0.6pt, decorate, decoration={brace, amplitude=5pt}] (7.2,-0.84) -- (11.68,-0.84);
\node[anchor=south, font=\fontfamily{phv}\selectfont\footnotesize\bfseries] at (9.44,-0.66) {$\times$32};
\node[blk, text width=3.6cm] (cb) at (9.3,-2.03) {Coupling bank, $2048\times64$\,b\\{\scriptsize 4 copies, 8 reads per cycle}};
\node[blk, fill=figred, text width=3.6cm] (pc) at (9.3,-3.13) {Population-count adder\\{\scriptsize $\sum_e 2\sigma_e J_{x_e y}$, no multipliers}};
\node[blk, text width=3.6cm] (fd) at (9.3,-4.23) {64 fields $h_y$, $y=b+32j$\\{\scriptsize 12 bit each}};
\node[blk, text width=3.6cm] (ln) at (9.3,-5.33) {8 decision lanes\\{\scriptsize PRNG, threshold, linear flag}};
\draw[bus] (cb) -- node[right, lab]{8 rows} (pc);
\draw[bus] (pc) -- (fd);
\draw[bus] (fd) -- (ln);
\draw[bus] (ld.east) -- ++(0.35,0) |- (cb.west);
\node[lab, anchor=east] at (6.7,-2.2) {$J$};
\draw[bus] (tab.east) -- ++(0.35,0) |- (ln.west);
\node[lab, anchor=east] at (6.7,-4.75) {constants};
\draw[bus] (7.2,-5.65) -- (sc.east|-0,-5.65);
\node[blk, fill=figyellow, text width=2.45cm, minimum height=1.3cm] (ex) at (13.55,-3.6) {8 flip extractors\\{\scriptsize priority encoders,\\up to 8 flips per cycle}};
\node[blk, fill=figyellow, text width=2.45cm] (nl) at (13.55,-5.3) {$n_{\mathrm{lin}}$ counter\\{\scriptsize population count $\to c(t)$}};
\draw[bus] (11.19,-5.15) -- (11.85,-5.15) -- (11.85,-3.9) -- (ex.west|-0,-3.9);
\node[lab, anchor=west] at (11.95,-4.55) {flip bits};
\draw[bus] (11.19,-5.48) -- (nl.west|-0,-5.48);
\draw[bus] (nl.south) -- (13.55,-6.2) -- (10.3,-6.2) -- (10.3,-5.74);
\node[lab, anchor=north] at (11.95,-6.24) {$c(t)$ for the correction};
\draw[bc] (ex.north) -- (13.55,-2.03) -- (cb.east);
\node[lab, anchor=south] at (12.65,-2.0) {8 flips/cycle\\$(x_e,\sigma_e)$ to\\all 32 groups};
\end{tikzpicture}}
\caption{(a)~Twelve engines, four per SLR, share the NoC and HBM of the AVED shell~\cite{amd_aved}, which connects to the host through PCIe. (b)~One engine. Each of the 32 groups (one shown) owns a coupling bank, 64 fields $h_y$ with $y=b+32j$ and 8 decision lanes, so coupling reads and field updates stay in the group. Eight extractors drain the flip bits while later rounds are still deciding and broadcast up to eight flipped spins $(x_e,\sigma_e)$ per cycle to all groups, which add $2\sigma_e J_{x_e y}$ to their fields by population counts. The count $n_{\mathrm{lin}}$ of linear decisions gives the coefficient $c(t)$ of the correction. PRNG: pseudorandom number generator.}
\Description{Left: a V80 device with three super logic regions, each holding four numbered engines, and a NoC, HBM and PCIe bar; engine 7 is highlighted and connected by zoom lines to the right panel. Right: one engine, with a front end, loader, step tables and score unit on the left; a stack of 32 group cards in the middle, the front card showing a coupling bank, a population-count adder, 64 fields and 8 decision lanes connected top to bottom; on the right, eight flip extractors that broadcast flipped spins back to the coupling banks of all groups, and an n_lin counter that returns the correction coefficient to the lanes.}
\label{fig:arch}
\vspace{-0.6em}%
\end{figure*}


\subsection{Data layout and decisions}\label{subsec:layout}
$J_{ij}\in\{\pm1\}$ is stored as one bit per entry. The padded $2048\times2048$ matrix is split by columns into 32 banks. Word $x$ of bank $b$ holds the 64 couplings $J_{x,\,b+32j}$ ($j=0..63$). Each bank exists four times, two copies in block random-access memory (BRAM) and two in UltraRAM (URAM), and each copy has two ports, which gives eight row reads per bank per cycle (16\,Mbit per engine). Group $b$ owns the 64 fields $h_y$ ($y=b+32j$, 12-bit), their spins, and eight decision lanes. Lane $l$ decides spin $l+256k$ in round $k$, so a step takes eight rounds of 256 parallel decisions.

Each lane has its own xoshiro128** generator~\cite{blackman2021scrambled}. At the start of a trial the generators are seeded through a 256-stage shift chain from one pipelined splitmix64 unit~\cite{steele2014splitmix}, which avoids a seed broadcast. A lane evaluates $z$ of Equation~\eqref{eq:onsager} in Q16.16 fixed point (16 integer and 16 fractional bits) and compares it with the threshold $((r-2^{15})\cdot4T_t)\gg16$, where $r$ is a 16-bit random number and $\gg$ denotes an arithmetic right shift. This needs one digital signal processing (DSP) multiplier per lane. The lane also reports whether the decision is linear ($|z|<2T_t$). Both comparisons are split into exact 16-bit halves to shorten the carry chains. The per-step correction is $(n_{\mathrm{lin}}(t{-}1)\cdot\gamma_t)\gg8$, where $\gamma_t=\lambda_t/(2T_{t-1})$ is stored in Q8.24 format (8 integer and 24 fractional bits). A second table entry holds the constant term of Onsager-$\kappa T$ or of TEC's coupling added to SCA~\cite{du2026tec}.

\subsection{Flip extraction and field update}\label{subsec:update}
Eight extractors each drain one 32-bit word of flip bits per round, eight lanes from each of four groups. A priority encoder emits one flipped spin $x_e$ and its new sign $\sigma_e$ per cycle per extractor. Extraction of round $k$ starts while rounds $k{+}1..7$ are still deciding. Each cycle the up to eight flips are broadcast to all groups. Every group reads the eight coupling words $J_{x_e,\cdot}$, one per bank port, and updates each of its 64 fields by $\sum_e 2\sigma_eJ_{x_e y}$. That sum is an 8-input population count followed by a 6-bit addition, so 16{,}000 coupling terms are applied per cycle without a single multiplier. Field updates become visible six cycles after issue; the next round reads its fields two cycles after the extractors drain, the shortest bit-exact latency.

The resulting cost of one trial in clock cycles, fitted on RTL simulation of one trial for each of six protocol configurations, is
\begin{equation}
  \text{cycles}=\cycFlip\cdot\text{flips}+\cycStep\cdot S+886 ,
  \label{eq:cycles}
\end{equation}
where flips is the total number of spin flips in the trial and $S$ is the number of steps. Each step costs \cycStep{} cycles to issue its eight decision rounds and drain the pipeline. Each flip adds \cycFlip{} cycles, slightly more than the ideal $1/8$, because a step waits for the busiest of the eight extractors. The constant covers the per-trial work outside the steps: seeding the generators, drawing the random initial state and building its fields, computing $\sum_is_ih_i$ and the cut in a pipelined adder tree, and streaming the result out. For example, a trial with $S=280$ and about 47{,}400 flips takes about 12{,}700 cycles in the model and 12{,}371 on the board on average (49\,\textmu s at 250\,MHz).
On the board, the fastest schedules average 38--44 cycles per step (153--177\,ns). Eight extractors and four coupling copies instead of four and two cut the cycles per trial by 28--32\% for 15\% more lookup tables (LUTs) and identical results (Table~\ref{tab:evo}).

\subsection{Front end and multi-engine system}\label{subsec:system}
An HLS front end gives each engine two AXI4 masters into the high-bandwidth memory (HBM) through the network on chip (NoC): a 256-bit master for couplings and tables and a second master for results. Couplings and tables stay resident across launches, so after the first launch a launch moves only a two-beat header.

The system places twelve engines, four per super logic region (SLR), with one placement block per SLR. The engines connect through two SmartConnects to the NoC of the Alveo Versal Example Design (AVED) shell~\cite{amd_aved}, which also provides the processing system, the PCIe interface and HBM. The kernel clock is 250\,MHz. Timing closure requires six changes: (1)~replicated clock enables and per-step constants, (2)~the split comparisons, (3)~a registered coupling-address multiplexer, (4)~a pipelined score tree, (5)~an input skid buffer, and (6)~post-route physical optimization.
These move the stand-alone four-extractor engine (v4) from $-0.78$\,ns to $+0.32$\,ns of slack at 300\,MHz; the eight-extractor engine alone meets 275\,MHz.

\subsection{Verification}\label{subsec:verif}
The engine reproduces a C++ reference model bit for bit: spins, cut, flip count and the per-step $n_{\mathrm{lin}}$ trace. This holds in RTL simulation for every protocol configuration. On the board, each image must first pass an exact gate. Seeded trials on every engine, with per-step traces, are compared bit for bit against the reference (60/60 on each 12-engine version). Because trial outcomes depend only on the seed and the trial index, images must also agree with each other. Versions v4--v6 (Table~\ref{tab:evo}) return identical results for all 41{,}040 common trials, and v7 for all 4{,}164 trials it shares with v6.

\subsection{Multi-bit couplings and bias}\label{subsec:multibit}
Sparse or weighted graphs need couplings beyond $\pm1$, and encodings with linear terms need a per-spin bias~\cite{lucas2014ising}. The multi-bit engine (v7 in Table~\ref{tab:evo}) stores $K$-bit couplings $J_{xy}\in\{-M,\dots,M\}$, $M=2^{K-1}-1$, as a sign plane and $K{-}1$ magnitude planes in the bank layout of Section~\ref{subsec:layout}. With $\tau_e=\sigma_e\operatorname{sign}(J_{x_ey})$ and $m_{be}$ bit $b$ of $|J_{x_ey}|$, the field update $\sum_e2\sigma_eJ_{x_ey}$ of Section~\ref{subsec:update} becomes $2\sum_b2^b\bigl(2\,\#\{e:m_{be}\wedge\tau_e{=}{+}1\}-\#\{e:m_{be}\}\bigr)$: two 8-input population counts per magnitude plane and still no multiplier. Each coupling is stored once, in a single-port bank per extractor, because extractor $e$ applies only flips of spins $x$ with $x\bmod8=e$. At $K=2$ this halves the stored bits of v6 (8 instead of 16\,Mbit per engine), but the store is limited by read ports rather than capacity: the engine keeps the BRAM of v6 and doubles its URAM. A bias $b_i$ enters only the initial fields, $h_i(0)=b_i+\sum_jJ_{ij}s_j(0)$; the incremental update preserves it, so it costs no cycles per step, and a runtime problem size $n\leq2{,}048$ masks unused spins. Twelve engines with $K=2$ close timing at 250\,MHz with 89\% of the LUTs and 81\% of the URAM. With $K=4$, the memory allows six engines, which fail routing at 250\,MHz but close timing at 225\,MHz (86\% of the LUTs). The $K=2$ engine is bit-exact with the reference model in simulation and on the board, including runtime sizes of 800--2{,}000 spins, and on K2000 it returns the same trials as v6 in the same cycles. With bias, the six $K=4$ engines are bit-exact on the board in 114 of 114 trials, including GP and TSP encodings, as is the GPU kernel (Section~\ref{subsec:gpu}; 6{,}308 trials).
